\documentclass[9pt,twocolumn,twoside]{pnas-new}

\templatetype{pnasresearcharticle}

% Compatibility shims for 2025 template commands not in local pnas-new.cls
\providecommand{\articletype}[1]{}
\providecommand{\Firstpage}{}
\providecommand{\Endparasplit}{}
\providecommand{\dataavail}[1]{\section*{Data, Materials, and Code Availability}#1}
\providecommand{\bibsplit}[1][]{}
\setboolean{displaywatermark}{false}

\articletype{RESEARCH ARTICLE}

\begin{document}

\title{Green Favoritism: Leader Birth Regions Receive Less Pollution-Intensive Growth}

\author[a,b]{Richard Schulz}
\author[a]{Muhammad Majid Altaf}
\author[a]{Ruby Wang}

\affil[a]{University of California, Berkeley, Berkeley, CA 94720, USA}
\affil[b]{ETH Zurich, 8092 Zurich, Switzerland}

\leadauthor{Schulz}

\significancestatement{Birth regions of national leaders emit more satellite-measured nightlight (a standard proxy for local economic activity) while their leader is in office, but without a proportional rise in measured air pollution. In a harmonized 2005--2017 panel, the nightlights coefficient is 0.049 ($p = 0.001$), NO$_2$ intensity falls by 0.039 ($p = 0.044$), and PM$_{2.5}$ intensity by 0.044 ($p = 0.004$). This is a relative shift in the pollution-to-nightlights ratio, not direct evidence of abatement. The pattern is substantially stronger in Sub-Saharan Africa and in lower-democracy countries, suggestive of a composition channel in which leaders direct large public works toward birth regions, though the satellite design cannot separate composition from selective enforcement or sorting.}

\authorcontributions{R.S. conceived the study, built and integrated the main data pipeline, ran the core empirical analysis, and drafted the manuscript. M.M.A. contributed to the NO$_2$ analysis. R.W. contributed to manuscript review and proofreading.}
\authordeclaration{The authors declare no competing interests.}
\correspondingauthor{To whom correspondence should be addressed. E-mail: rschul@ethz.ch}
\keywords{regional favoritism $|$ political economy $|$ nighttime lights $|$ air pollution $|$ green favoritism}
\acknow{We thank Professor Joshua Blumenstock for feedback. All errors are our own.}

\begin{abstract}
We replicate the birth-region nightlights favoritism documented by Bora (2025) and D\"{u}ben, Hodler, and Raschky (2026) and extend the result to air pollution. We define \emph{green favoritism} as a negative birth-region coefficient on pollution intensity, defined as $\log(\text{Pollution}+0.01) - \log(\text{Nightlights}+0.01)$. A negative coefficient means birth-region growth is more nightlight-intensive relative to measured air pollution than non-birth-region growth in the same country-year; it is a relative shift in the pollution-to-nightlights ratio, not direct evidence of abatement. Using a harmonized DMSP-to-VIIRS panel for 2005--2017 (622,868 district-year observations), the birth-region coefficient on log nightlights is 0.049 (SE = 0.015). NO$_2$ intensity falls by 0.039 (SE = 0.019) and PM$_{2.5}$ intensity by 0.044 (SE = 0.015) at leader entry. An event study on the same 2005--2017 samples shows increasingly positive nightlights and increasingly negative pollution-intensity coefficients over leaders' time in office, though the post-exit coefficients remain elevated. The effect is substantially stronger in Sub-Saharan Africa, where the nightlights coefficient is roughly three times the non-SSA estimate, although the non-SSA coefficient is itself positive and significant. The pollution-intensity effect is also stronger in lower-democracy countries, suggestive of politically discretionary channels (composition or selective enforcement) that the satellite data cannot separately identify.
\end{abstract}

\maketitle
\thispagestyle{firststyle}
\ifthenelse{\boolean{shortarticle}}{\ifthenelse{\boolean{singlecolumn}}{\abscontentformatted}{\abscontent}}{}

\Firstpage

Political leaders direct resources toward their birth regions. Cross-country panels of subnational nightlights have documented this pattern repeatedly: districts that are home to the sitting national leader grow brighter relative to comparable districts in the same country-year \cite{HodlerRaschky2014, Bora2025, DubenHodlerRaschky2026}. A natural follow-up question, however, has not been examined: \emph{is birth-region growth more or less pollution-intensive than activity elsewhere?} If leaders channel public goods and infrastructure toward home districts, pollution should rise less than economic activity. If they route industrial subsidies or manufacturing instead, pollution should rise proportionally or more.

We find that birth-region growth is, in this relative sense, less pollution-intensive. Using a harmonized DMSP-to-VIIRS panel for 2005--2017, we first replicate the nightlights result with a birth-region coefficient of 0.049 (SE = 0.015, $p = 0.001$). Both NO$_2$ intensity and PM$_{2.5}$ intensity---defined as $\log(\text{Pollution}+0.01) - \log(\text{Nightlights}+0.01)$---are significantly negative: $-0.039$ (SE = 0.019) and $-0.044$ (SE = 0.015). Birth-region growth is more nightlight-intensive relative to measured air pollution than activity in the rest of the same country-year, consistent with a lower pollution-to-nightlights ratio rather than direct evidence of abatement. We term this pattern \emph{green favoritism}.

An event study on the same 2005--2017 headline samples provides a dynamic check on the main specification. The pre-entry coefficients are modest and statistically insignificant for all three outcomes, while the in-office coefficients become more positive for nightlights and more negative for both pollution-intensity measures over leaders' tenure. Post-exit coefficients remain elevated rather than snapping back immediately, so we read the figure as supportive of gradual in-office buildup but not as a sharp causal timing test on its own.

Hodler and Raschky (2014) establish birth-region favoritism using DMSP nightlights. Bora (2025) and D\"{u}ben, Hodler, and Raschky (2026) confirm the result with updated birthplace data and extended samples. De Luca et al.\ (2018) extend the favoritism result to ethnic homelands \cite{DeLuca2018}, and Bomprezzi et al.\ (2024) document a parallel pattern for the spousal birthplaces of leaders \cite{Bomprezzi2024Spouses}. The literature on regional favoritism is, in this sense, well populated on the side of \emph{how much} activity birth regions receive, but largely silent on the \emph{kind} of activity they receive.

A separate literature studies the political economy of pollution. Cross-country work links environmental outcomes to trade openness, factor endowments, and the strength of regulatory institutions \cite{CopelandTaylor2004}. Our paper sits at the intersection of this literature and the favoritism literature: the favoritism literature has examined how much activity birth regions receive but not what kind, while the pollution literature has examined cross-country and within-country drivers of emissions but not the role of leader birthplace. The contribution here is to show that political connection shapes not only the level of activity but its measured pollution-to-nightlights ratio. The pollution-intensity outcome uses independent satellite data on air quality, but it is not a separate identifying margin: under the same regressor and fixed-effects structure the intensity coefficient is mechanically $\beta_{\text{pollution}} - \beta_{\text{NTL}}$, so the substantive value of the intensity outcome is in characterizing the kind of growth birth regions receive rather than in providing a new source of identification.

This paper documents green favoritism, a pollution-intensity pattern characterized using independent satellite data on air quality. The effect is substantially stronger in Sub-Saharan Africa, where the nightlights coefficient is roughly three times the non-SSA estimate, though the non-SSA estimate is itself positive and significant.

Section 1 describes the data. Section 2 lays out the empirical strategy. Section 3 presents the replication. Section 4 documents green favoritism. Section 5 turns to the event study. Section 6 examines geographic heterogeneity, and Section 7 discusses mechanisms.\Endparasplit

\section{Data}

The analysis combines five core inputs. First, the Political Leaders' Affiliation Database (PLAD) provides leader identity, spell timing, and subnational birthplace information used to code treatment \cite{Bomprezzi2025PLAD}. Second, nightlights provide the economic-activity proxy. For replication we use a DMSP ADM2 panel for 1992--2013. For the main analysis we use a harmonized DMSP-to-VIIRS panel for 2005--2017. DMSP records luminance as an integer Digital Number (DN) from 0 to 63, with known saturation near the ceiling \cite{Gibson2021}. Our zonal statistics pipeline treats zero-valued pixels as valid dark observations rather than missing data, preserving the extensive margin of electrification. Third, we use ACAG surface NO$_2$ data for 2005--2017 \cite{Cooper2022NO2}. Fourth, we use ACAG SatPM PM$_{2.5}$ as a complementary pollution measure. Fifth, GADM ADM2 boundaries provide the common subnational geography \cite{GADM2022}.

Raster extraction follows a uniform pipeline across all satellite products. Annual rasters are aggregated to ADM2 polygons by computing the mean of all valid pixels intersecting each boundary polygon (zonal statistics). Zero-valued pixels (districts with no detectable nightlights or surface concentrations) are treated as valid observations rather than missing data, which preserves the extensive margin of electrification and retains variation across dark, low-pollution regions. Only pixels designated \texttt{nodata} by the source product (ocean, ice, and out-of-bounds cells) are excluded from the mean. Outcome variables enter regressions as $\log(Y_{ict} + 0.01)$; the constant shift avoids undefined logarithms for zero-valued observations and maps dark districts to approximately $-4.6$. Results are robust to alternative shift values of 0.001 and 1.

Table~\ref{tab:sumstats} reports summary statistics for the 2005--2017 analysis sample.

\begin{table}[t]
\centering
\caption{Summary statistics, 2005--2017 analysis sample}
\small
\begin{tabular*}{\columnwidth}{@{\extracolsep{\fill}}lrrr}
Variable & Mean & SD & $N$ \\
\midrule
Log nightlights (harmonized) & 0.508 & 2.351 & 622{,}868 \\
NO$_2$ (ppb) & 0.972 & 2.079 & 585{,}141 \\
PM$_{2.5}$ ($\mu$g/m$^3$) & 18.051 & 11.472 & 620{,}805 \\
Birth-region indicator & 0.003 & 0.057 & 622{,}868 \\
\bottomrule
\end{tabular*}
\addtabletext{Unit of observation is ADM2 district-year. The birth-region indicator equals one when the district is the birth region of the sitting national leader. Observations differ across variables due to differences in satellite coverage.}
\label{tab:sumstats}
\end{table}

\section{Empirical Strategy}

Treatment is defined at the district-year level. A district is treated when it is the birth region of the national leader in office in that year. To ensure a unique treatment assignment in transition years, we truncate the earlier of any two consecutive leader spells that overlap within a country.

The core estimating equation uses ADM2 fixed effects and country-year fixed effects:
\[
\log(Y_{i c t} + 0.01) = \beta \, \text{BirthRegion}_{i c t} + \alpha_i + \gamma_{c t} + \varepsilon_{i c t},
\]
where $i$ indexes ADM2 regions, $c$ countries, and $t$ years. The ADM2 fixed effects $\alpha_i$ absorb all time-invariant district characteristics: geography, historical institutions, and long-run income. The country-year fixed effects $\gamma_{ct}$ absorb national policy shocks, macroeconomic cycles, and any aggregate trend in pollution or nightlights common to all districts of a country in a given year. The coefficient $\beta$ is therefore identified by within-country-year variation: whether a district that becomes the current leader's birth region experiences differential growth relative to other districts in the same country at the same time.

Pollution intensity is $\log(\text{Pollution}+0.01) - \log(\text{Nightlights}+0.01)$, measuring emission output per unit of proxied economic activity. Because it uses the same regressor and fixed effects as the component regressions, its coefficient is mechanically $\beta_{\text{pollution}} - \beta_{\text{NTL}}$. A negative intensity coefficient means pollution rises less than nightlights in birth regions.

The main threat to identification is pre-existing trends. Leaders may disproportionately originate from districts already on different growth trajectories, so that the post-treatment levels reflect pre-existing conditions rather than causal favoritism. This is the selection concern raised by Hodler and Raschky (2014), and the event study in Section 5 tests it directly by examining pre-period coefficients. The pollution-intensity outcome carries an additional advantage: because it differences out the nightlights component, common drivers of both nightlights and pollution (including most pre-existing economic trends) cancel. This makes the green favoritism coefficient harder to generate spuriously than the nightlights coefficient alone.

Standard errors are clustered by leader spell in the baseline, following the convention in recent replication work \cite{Bora2025}.

\section{Replication}

Table~\ref{tab:replication} presents a replication ladder for the birth-region nightlights effect across successive sample windows and sensor generations. In the long DMSP 1992--2013 window the estimate is 0.043 ($p = 0.002$); in the overlap DMSP 2005--2013 window, 0.033 ($p = 0.019$); in the harmonized 2005--2017 panel, 0.049 ($p = 0.001$). The estimates are of similar magnitude and all statistically significant across sensor generations, consistent with Bora (2025) and D\"{u}ben, Hodler, and Raschky (2026). The harmonized 2005--2017 specification in row (3) is the one we extend to pollution outcomes in Section 4.

\begin{table}[t]
\centering
\caption{Replication ladder: birth-region nightlights effect across sensor generations}
\small
\begin{tabular*}{\columnwidth}{@{\extracolsep{\fill}}llrrr}
Step & Outcome & Coef. & SE & $p$ \\
\midrule
(1) DMSP 1992--13 & Log NTL & 0.043 & 0.014 & 0.002 \\
(2) DMSP 2005--13 & Log NTL & 0.033 & 0.014 & 0.019 \\
(3) Harm.\ 2005--17 & Log NTL & 0.049 & 0.015 & 0.001 \\
\bottomrule
\end{tabular*}
\addtabletext{Lag-0 birth-region coefficients on log nightlights across sensor generations ($N = 1{,}054{,}103$; $431{,}208$; $622{,}868$). All rows: ADM2 FE, country-year FE, spell-clustered SE. Row (3) is the specification carried forward to the pollution outcomes in Table~\ref{tab:green-fav}.}
\label{tab:replication}
\end{table}

\section{Green Favoritism}

Table~\ref{tab:green-fav} presents the main result. In the harmonized 2005--2017 panel, the birth-region coefficient on log nightlights is 0.049 (SE = 0.015, $p = 0.001$), directly replicating the result in Table~\ref{tab:replication}. The new result is the pollution side: the birth-region coefficient on NO$_2$ intensity is $-0.039$ (SE = 0.019, $p = 0.044$) and on PM$_{2.5}$ intensity is $-0.044$ (SE = 0.015, $p = 0.004$). Both pollution-intensity coefficients are negative and statistically significant. Birth regions receive economic activity that does not come with proportional air pollution.

\begin{table}[t]
\centering
\caption{Green favoritism: birth regions grow brighter and less pollution-intensive (2005--2017)}
\small
\begin{tabular*}{\columnwidth}{@{\extracolsep{\fill}}lrrr}
 & (1) & (2) & (3) \\
 & $\ln$ NTL & NO$_2$ int. & PM$_{2.5}$ int. \\
\midrule
Birth-region indicator & $0.049^{***}$ & $-0.039^{*}$ & $-0.044^{**}$ \\
 & $(0.015)$ & $(0.019)$ & $(0.015)$ \\
\midrule
$N$ & 622,868 & 585,141 & 620,805 \\
Region FE & Yes & Yes & Yes \\
Country-year FE & Yes & Yes & Yes \\
SE clustering & Spell & Spell & Spell \\
\bottomrule
\end{tabular*}
\addtabletext{Pollution intensity is $\log(\text{Pollution}+0.01) - \log(\text{Nightlights}+0.01)$. Negative intensity coefficients mean pollution rises less than nightlights in birth regions. Observations differ across columns because NO$_2$ and PM$_{2.5}$ have different spatial coverage. $^{*}p<0.05$; $^{**}p<0.01$; $^{***}p<0.001$. Standard errors clustered by leader spell.}
\label{tab:green-fav}
\end{table}

The green favoritism result does not require that pollution falls in absolute terms. The lag-0 NO$_2$ level coefficient is small and insignificant (0.006, $p = 0.585$); the lag-0 PM$_{2.5}$ level coefficient is similarly small and not significantly different from zero. The negative intensity is driven by a large positive nightlights coefficient combined with a near-zero pollution response: birth regions grow brighter without growing proportionally dirtier. The intensity coefficients reported in Table~\ref{tab:green-fav} are estimated directly from the intensity outcome rather than constructed by subtraction; because the level and intensity regressions are run on slightly different effective samples (NO$_2$ and PM$_{2.5}$ have different spatial coverage from nightlights), the intensity coefficient is not exactly equal to the difference of the level coefficients reported in the text.

The two pollution indicators are complementary rather than interchangeable. The NO$_2$-intensity estimate is sensitive to country-level clustering and loses statistical significance under that more conservative scheme (Table~\ref{tab:cluster-robustness}), while the PM$_{2.5}$-intensity estimate is more stable. In the event study, both pollution-intensity measures show modest pre-entry coefficients and increasingly negative in-office coefficients, but the post-exit coefficients remain elevated (Section~5). PM$_{2.5}$ captures a broader mix of emission sources including transportation and combustion, while NO$_2$ is more specific to traffic and industrial combustion.

Table~\ref{tab:lag-table} shows that the effects persist across lags. The nightlights effect peaks at lag 1 (+0.063, $p < 0.001$), consistent with a one-year delay between fiscal allocation and measurable luminosity. Both pollution-intensity effects remain negative and statistically significant through lag 2, suggesting some persistence in the lower pollution-to-nightlights ratio in the full-panel lag specification.

\begin{table}[t]
\centering
\caption{Green favoritism across lags (2005--2017)}
\small
\begin{tabular*}{\columnwidth}{@{\extracolsep{\fill}}lrrr}
 & (1) & (2) & (3) \\
 & $\ln$ NTL & NO$_2$ int. & PM$_{2.5}$ int. \\
\midrule
Lag 0 & $0.049^{***}$ & $-0.039^{*}$ & $-0.044^{**}$ \\
      & $(0.015)$     & $(0.019)$    & $(0.015)$ \\
Lag 1 & $0.063^{***}$ & $-0.081^{***}$ & $-0.059^{***}$ \\
      & $(0.015)$     & $(0.019)$      & $(0.015)$ \\
Lag 2 & $0.042^{**}$  & $-0.070^{***}$ & $-0.039^{*}$ \\
      & $(0.016)$     & $(0.021)$      & $(0.016)$ \\
\midrule
$N$              & 622,868 & 585,141 & 620,805 \\
Region FE        & Yes & Yes & Yes \\
Country-year FE  & Yes & Yes & Yes \\
SE clustering    & Spell & Spell & Spell \\
\bottomrule
\end{tabular*}
\addtabletext{Each cell is a separate regression. Lag $k$ uses the birth-region indicator shifted $k$ years forward. Pollution intensity is $\log(\text{Pollution}+0.01) - \log(\text{Nightlights}+0.01)$. All rows use harmonized 2005--2017 panel, ADM2 fixed effects, country-year fixed effects, and standard errors clustered by leader spell. $^{*}p<0.05$; $^{**}p<0.01$; $^{***}p<0.001$.}
\label{tab:lag-table}
\end{table}

\section{Event Study}

Figure~\ref{fig:event-study} shows an HR-style event study for log nightlights, NO$_2$ intensity, and PM$_{2.5}$ intensity on the same 2005--2017 samples used in the headline regressions. The figure comes from a single regression per outcome with dummies for the three years before leader entry, individual in-office years 0--5, a pooled in-office 6+ bin, and the first three post-exit bins, with $t=-1$ omitted. The solid red line marks entry and the dashed red line marks the transition from in-office to post-exit bins.

The resulting profiles are consistent with the static estimates but do not show a textbook jump at entry. For nightlights, the pre-entry coefficients are small, and the in-office coefficients become increasingly positive over tenure, reaching 0.139 at year 5 and 0.159 in the pooled 6+ bin. For NO$_2$ and PM$_{2.5}$ intensity, the in-office coefficients are negative and become larger in magnitude over time: NO$_2$ reaches $-0.201$ at year 5 and $-0.198$ in 6+, while PM$_{2.5}$ reaches $-0.133$ at year 5 and $-0.152$ in 6+. Post-exit coefficients remain elevated rather than reverting immediately, so the event study supports a gradual buildup in birth-region activity and relatively cleaner composition while also cautioning against a sharp causal timing interpretation. PM$_{2.5}$ remains the more robust pollution indicator under conservative country-level inference (Table~\ref{tab:cluster-robustness}), while NO$_2$ is the more source-specific measure of traffic and industrial combustion.

\begin{figure*}[t]
\centering
\includegraphics[width=\textwidth]{figures/event_study.pdf}
\caption{HR-style event study on the 2005--2017 headline samples. Three panels show log nightlights, NO$_2$ intensity, and PM$_{2.5}$ intensity. Coefficients are relative to $t=-1$ and come from a single regression per outcome with pre-entry bins ($-3,-2$), individual in-office bins ($0$ through $5$), a pooled in-office 6+ bin, and post-exit bins ($+1,+2,+3+$). The dark blue line plots the coefficient estimates, and the light blue lines indicate the upper and lower 95\% confidence interval bounds. The solid red line marks entry, the dashed red line marks the transition from in-office to post-exit bins, and the dashed horizontal line marks the main static coefficient for that outcome.}
\label{fig:event-study}
\end{figure*}

The persistence of the post-exit coefficients suggests either slow-moving project dynamics or lingering differences in district trajectories after leaders leave office. Distinguishing those possibilities is beyond the scope of the satellite panel and remains an open question.

\section{Geographic Heterogeneity}

The birth-region nightlights effect is substantially stronger in Sub-Saharan Africa than elsewhere, though it is not absent outside SSA. Table~\ref{tab:ssa-decomp} decomposes the 2005--2017 nightlights effect by region. The SSA coefficient is 0.109 ($p = 0.017$, $N = 56{,}589$), roughly three times the non-SSA estimate of 0.033 ($p = 0.026$, $N = 566{,}279$) and about 2.2 times the pooled estimate of 0.049. The non-SSA coefficient is itself positive and statistically significant; SSA is the more pronounced case rather than the only one. The SSA pollution-intensity estimates are also more negative: NO$_2$ intensity at $-0.106$ (SE = 0.051, $p = 0.042$) and PM$_{2.5}$ intensity at $-0.112$ (SE = 0.044, $p = 0.013$). We treat the SSA heterogeneity results cautiously, since both satellite measurement and ADM2 aggregation may differ across SSA and non-SSA settings.

\begin{table}[t]
\centering
\caption{Geographic decomposition: 2005--2017 nightlights}
\small
\begin{tabular*}{\columnwidth}{@{\extracolsep{\fill}}lrrrr}
Sample & Coef. & SE & $p$-value & $N$ \\
\midrule
Pooled & 0.049 & 0.015 & 0.001 & 622{,}868 \\
Sub-Saharan Africa & 0.109 & 0.046 & 0.017 & 56{,}589 \\
Non-SSA & 0.033 & 0.015 & 0.026 & 566{,}279 \\
Africa & 0.101 & 0.040 & 0.012 & 85{,}063 \\
Low NO$_2$ dispersion & 0.108 & 0.033 & 0.001 & 92{,}499 \\
\bottomrule
\end{tabular*}
\addtabletext{Lag-0 birth-region coefficients on log nightlights, harmonized 2005--2017. All rows use ADM2 fixed effects, country-year fixed effects, and leader-spell clustered standard errors. Low-NO$_2$-dispersion countries are defined as countries below the median cross-regional standard deviation of NO$_2$ in 2005--2007.}
\label{tab:ssa-decomp}
\end{table}

The SSA concentration is consistent across replication windows: the SSA coefficient is 0.125 in DMSP 1992--2013, 0.116 in DMSP 2005--2013, and 0.109 in harmonized 2005--2017, in each case substantially larger than the corresponding non-SSA estimate.

\section{Mechanism}

The green favoritism pattern is consistent with leaders channeling investment with low emission intensity toward birth regions, such as large public works (ports, airports, government facilities) that generate nightlights without proportional industrial emissions. The aggregate satellite design used here cannot directly identify project composition, and the case of Sri Lanka under President Mahinda Rajapaksa (2005--2015) is offered as illustrative rather than as validation of the mechanism. Rajapaksa was born in Hambantota District; during his presidency, the government constructed the Hambantota International Port, Mattala Rajapaksa International Airport, and a major convention center in that district, infrastructure projects that substantially increased local luminosity without industrial activity. Whether this pattern generalizes across countries requires project-level investment data that we do not have.

Three channels could generate negative pollution intensity in birth regions. The first is a \emph{composition channel}: leaders direct low-emission investment (ports, airports, roads, government facilities) toward birth regions instead of heavy manufacturing or industrial plants. The second is an \emph{enforcement channel}: birth-region firms may face stricter environmental enforcement while the leader is in office, reducing emissions per unit of output. The third is a \emph{sorting channel}: birth regions may attract already-clean sectors (services, tourism, knowledge industries) rather than manufacturing. The satellite data used here cannot distinguish these channels, and the democracy split below is consistent with both composition and selective enforcement: low-democracy environments may permit either more discretionary project allocation or more discretionary enforcement, and the aggregate design cannot tell them apart.

\subsection*{Democracy and accountability}

Table~\ref{tab:democracy} splits the sample at the median country-mean V-Dem electoral democracy score (0.499) over 2005--2017. The nightlights favoritism effect is present in both groups, but the green favoritism pattern is concentrated in low-democracy countries. The NO$_2$ intensity coefficient in high-democracy countries is small and insignificant ($-0.021$, $p = 0.35$), while in low-democracy countries it is $-0.081$ ($p = 0.031$). The PM$_{2.5}$ intensity result is significant in both groups but roughly twice as large in low-democracy contexts.

\begin{table}[t]
\centering
\caption{Democracy heterogeneity (2005--2017)}
\small
\begin{tabular*}{\columnwidth}{@{\extracolsep{\fill}}lrrr}
 & (1) & (2) & (3) \\
 & $\ln$ NTL & NO$_2$ int. & PM$_{2.5}$ int. \\
\midrule
Low democracy & $0.066^{*}$ & $-0.081^{*}$ & $-0.062^{*}$ \\
              & $(0.029)$   & $(0.038)$    & $(0.029)$ \\
High democracy & $0.042^{*}$ & $-0.021$ & $-0.036^{*}$ \\
               & $(0.018)$   & $(0.023)$ & $(0.018)$ \\
\midrule
$N$ (low / high) & 203k / 408k & 185k / 391k & 202k / 407k \\
\bottomrule
\end{tabular*}
\addtabletext{Sample split at the median country-mean V-Dem electoral democracy score (0.499) over 2005--2017 \cite{VDem2025}. Pollution intensity is $\log(\text{Pollution}+0.01) - \log(\text{Nightlights}+0.01)$. Each cell is a separate regression with ADM2 fixed effects, country-year fixed effects, and standard errors clustered by leader spell. $^{*}p<0.05$.}
\label{tab:democracy}
\end{table}

This pattern is consistent with politically discretionary channels in settings with weaker accountability, including both project composition and selective enforcement. The aggregate satellite design cannot cleanly distinguish these mechanisms; both predict stronger green favoritism where leaders face less monitoring.

\subsection*{Robustness}

All main tables cluster standard errors by leader spell, following the convention in recent replication work \cite{Bora2025}. Table~\ref{tab:cluster-robustness} reports the same coefficients with country-level clustering as a more conservative robustness check. The nightlights coefficient in the harmonized window remains significant ($p = 0.013$). The PM$_{2.5}$ intensity result also survives ($p = 0.025$), while NO$_2$ intensity loses precision ($p = 0.119$). PM$_{2.5}$ thus provides the more conservative evidence for green favoritism under country-level inference.

\begin{table}[t]
\centering
\caption{Robustness to country-level clustering}
\small
\begin{tabular*}{\columnwidth}{@{\extracolsep{\fill}}lrrrr}
Spec. & Coef. & Spell SE & Ctry SE & Ctry $p$ \\
\midrule
NL, 1992--2013 & 0.043 & 0.014 & 0.018 & 0.016 \\
NL, 2005--2013 & 0.033 & 0.014 & 0.017 & 0.053 \\
NL, 2005--2017 & 0.049 & 0.015 & 0.020 & 0.013 \\
NO$_2$ int., 2005--2017 & $-0.039$ & 0.019 & 0.025 & 0.119 \\
PM$_{2.5}$ int., 2005--2017 & $-0.044$ & 0.015 & 0.020 & 0.025 \\
\bottomrule
\end{tabular*}
\addtabletext{NL denotes nightlights. Baseline standard errors cluster by leader spell; the robustness column clusters by country, following the original Hodler--Raschky convention.}
\label{tab:cluster-robustness}
\end{table}

\subsection*{Limitations}

Several caveats apply. First, satellite NO$_2$ and PM$_{2.5}$ are vertically integrated retrievals translated to surface concentrations by chemical transport models; they are noisy proxies for ground-level exposure and are known to perform less well over bright surfaces (deserts, snow) and over very dark, low-emission regions \cite{Cooper2022NO2}. Classical measurement error in a covariate would attenuate coefficients toward zero, but for log-transformed outcomes the bias direction is not generally signable, so we do not lean on attenuation as a defense of the pollution-intensity result. Second, nightlights are an imperfect proxy for economic activity, especially at low luminance and in agricultural regions where output need not translate into emitted light. The harmonized DMSP-to-VIIRS series mitigates the saturation issues of raw DMSP \cite{Gibson2021} but does not eliminate them. Third, birthplace coding in PLAD relies on biographical sources of varying quality; misclassification of a leader's birth region would generally attenuate the treatment indicator. Fourth, the PM$_{2.5}$ result is statistically more robust than the NO$_2$ result under country-level clustering, and we therefore treat PM$_{2.5}$ as the more conservative evidence for green favoritism, while NO$_2$ remains useful as the more source-specific pollutant. Fifth, the event-study profiles show gradual buildup rather than a sharp jump at entry, and the post-exit coefficients remain elevated, so the dynamic figure should be read as descriptive rather than as a sharp causal timing test. Sixth, because treatment timing is staggered and leadership spells turn on and off, the two-way fixed-effects estimates may be affected by recent concerns about forbidden comparisons in staggered DiD settings; robustness using modern staggered-adoption estimators (Callaway--Sant'Anna; de Chaisemartin--D'Haultfoeuille) is an important extension we leave for future work. Finally, the paper documents a pattern but cannot isolate the underlying channel. Composition, enforcement, and sorting all remain consistent with the aggregate evidence, and distinguishing them requires project-level or firm-level data not available at the global panel scale used here.

\section*{Discussion}

This paper asked whether birth-region growth is more or less pollution-intensive than activity elsewhere. We find that it is less so, in the relative sense that the pollution-to-nightlights ratio is lower in birth regions than in comparable districts of the same country-year. Birth-region favoritism in nightlights is a durable feature of the data, present across sensor generations, sample windows, and clustering schemes. The new result is that this favoritism is environmentally selective: birth regions grow brighter without a proportional rise in measured air pollution. Negative pollution-intensity coefficients, increasingly negative in-office event-study profiles, and the SSA concentration all point to birth-region growth that differs in composition---or in measured emissions per unit of luminosity---from growth elsewhere. The democracy split supplies one supporting strand of evidence: pollution-intensity favoritism is concentrated in lower-democracy countries, consistent with politically discretionary channels in settings with weaker accountability. The data do not let us cleanly isolate composition from selective enforcement or sectoral sorting; the aggregate pattern is suggestive of compositional differences in leader birth-region investment but does not by itself identify the channel.

The finding bears on environmental inequality, but only in a relative sense. The estimates suggest a shift in the pollution-to-nightlights ratio away from leader birth regions; the welfare and exposure implications require population-weighted analysis of pollution levels that we do not undertake here. These results should therefore be read as descriptive evidence of differential pollution-to-nightlights trajectories around leader birth regions, not as definitive evidence of abatement or of a uniquely identified mechanism. Despite these limitations, the paper adds an environmental dimension to the birth-region favoritism literature: political favoritism is visible not only in activity proxies but also in the relationship between activity and satellite-measured air pollution. Future work could test whether specific project categories---port and airport construction, electrification, road building---account for the lower-intensity pattern by linking nightlights and pollution outcomes to administrative investment records, and could revisit the staggered DiD inference using modern estimators robust to forbidden comparisons.

\matmethods{Data sources and panel construction are described in Section 1. The nightlights panels are built at ADM2 resolution by extracting annual zonal means from DMSP DN rasters (1992--2013) and harmonized DMSP-to-VIIRS rasters (2005--2017) using GADM level-2 boundaries. NO$_2$ and PM$_{2.5}$ are extracted from ACAG annual surface-concentration grids using the same boundaries. Leader birthplace districts are coded from PLAD and matched to GADM ADM2 polygons. The democracy split uses country-mean V-Dem electoral democracy scores (v2x\_polyarchy) over 2005--2017 \cite{VDem2025}, with countries split at the cross-country median (0.499). All regressions use \texttt{linearmodels} PanelOLS with absorbed fixed effects; standard errors are clustered by leader spell unless otherwise noted.}

\showmatmethods{}

\dataavail{Code and manuscript sources are available at \url{https://github.com/RichSchulz/political-leader-pollution}. External inputs include PLAD \cite{Bomprezzi2025PLAD}, GADM \cite{GADM2022}, ACAG NO$_2$ \cite{Cooper2022NO2} and PM$_{2.5}$, Li et al.\ harmonized nightlights, and DMSP annual composites obtained under the providers' access conditions.}

\showacknow{}

\bibsplit[3]

\bibliography{references}

\end{document}
